\subsection{Критерий компактности}
\begin{definition}[$\varepsilon$-сеть]\label{varepsilon-net}
	Пусть $M$~--- некоторое множество в метрическом пространстве $R$. Тогда множество $A$ из $R$ называется \textbf{$\varepsilon$-сетью} для $M$, если для любой точки $x \in M$ найдется хотя бы одна точка $a \in A$, что
	\[\rho(x, a) \leqslant \varepsilon.\]
\end{definition}

\begin{definition}[Ограниченность]
	Множество $M$ в метрическом пространстве $R$ называется \textbf{ограниченным}, если существует $B_\varepsilon(x_0)$ содержащий $M$.
\end{definition}

\begin{definition}[Вполне ограниченность, полная ограниченность]
	Множество $M$ в метрическом пространстве $R$ называется \textbf{вполне ограниченным}, если для него при любом $\varepsilon > 0$ существует \textbf{конечная $\varepsilon$-сеть}.
\end{definition}

\begin{remark}
	Если множество вполне ограничено, то его замыкание также вполне ограничено.
\end{remark}

\begin{remark}
	Из определения следует, что если само пространство вполне ограничено, то оно сепарабельно.
\end{remark}

\begin{remark}
	Из вполне ограниченности следует ограниченность.
\end{remark}

\begin{proof}
	Из вполне ограниченности ограниченность получается как объединение конечного числа ограниченных множеств.
\end{proof}

\begin{remark}
	В $n$-мерном евклидовом пространстве вполне ограниченность совпадает с обычной ограниченностью.
\end{remark}

\begin{theorem}[Критерий компактности]\label{thm3.2}
	Пусть $X$ --- метрическое пространство. Тогда следующие условия эквивалентны:
	\begin{enumerate}
		\item $X$ компактно;
		\item $X$ полно и вполне ограниченно;
		\item Из любой последовательности $\{x_n\} \subset X$ можно выделить сходящуюся подпоследовательность $\{x_{n_k}\}$ (т.е. $X$ — \textit{секвенциально компактно});
		\item Любое бесконечное множество $M \subset X$ имеет предельную точку.
	\end{enumerate}
\end{theorem}

\begin{proof}
	\begin{itemize}
		\item[\textbf{(1) $\Rightarrow$ (2)}]
		\begin{enumerate}
			\item \textbf{Вполне ограниченность.} Для любого $\varepsilon > 0$ рассмотрим покрытие $\{B(x, \varepsilon)\}_{x \in X}$. В силу компактности $X$, выделим конечное подпокрытие. Центры этих шаров образуют конечную $\varepsilon$-сеть.
			
			
			
			\item \textbf{Полнота.} Пусть $\{x_n\}$ — фундаментальная последовательность. Рассмотрим множества $F_n = \overline{\{x_k\}_{k \geqslant n}}$.
			\begin{itemize}
				\item Система $\{F_n\}$ — центрированная система замкнутых множеств (пересечение любого набора содержит хвост последовательности).
				\item По Теореме~\ref{thm3.1}, $\bigcap_{n} F_n \neq \varnothing$. Пусть $x_0 \in \bigcap F_n$.
				\item Так как диам($F_n$) $\to 0$ (в силу фундаментальности), то $x_0$ единственна и $x_n \to x_0$.
			\end{itemize}
		\end{enumerate}
		
		\item[\textbf{(2) $\Rightarrow$ (3)}] Пусть $\{x_n\} \subset X$.
		\begin{enumerate}
			\item \textbf{Построение.} Так как $X$ вполне ограничено, его можно покрыть конечным числом шаров радиуса $1$. В одном из них ($B_1$) содержится бесконечно много членов $\{x_n\}$.
			\item Далее покрываем $B_1$ шарами радиуса $1/2$. В одном из них ($B_2 \subset B_1$) снова бесконечно много членов.
			\item Продолжая процесс, получим вложенные шары $B_1 \supset B_2 \supset \dots$ с радиусами $r_k \to 0$, содержащие подпоследовательность $\{x_{n_k}\}$ (выбираем $x_{n_k} \in B_k$).
			\item \textbf{Сходимость.} $\{x_{n_k}\}$ фундаментальна (так как лежит в шарах стягивающегося радиуса). В силу полноты $X$, она сходится.
		\end{enumerate}
		
		\item[\textbf{(3) $\Rightarrow$ (1)}]
		\begin{enumerate}
			\item \textbf{Вполне ограниченность.} От противного: если конечной сети нет, то можно построить последовательность «растолканных» точек: $x_1$, затем $x_2$ на расстоянии $\rho > \varepsilon_0$ от первой, $x_3$ на расстоянии $\rho > \varepsilon_0$ от первых двух и т.д. У такой «редкой» последовательности нет сходящейся подпоследовательности. Противоречие.
			
			\item \textbf{Компактность.} От противного: пусть покрытие $\{G_\alpha\}$ не имеет конечного подпокрытия. 
			Будем уменьшать масштаб:
			\begin{itemize}
				\item Разделим всё пространство на конечное число мелких шаров. Хотя бы один из них ($B_1$) — «плохой» (не покрывается конечно).
				\item Разделим этот «плохой» шар на еще более мелкие. Снова найдем «плохой» подшар $B_2 \subset B_1$.
				\item Повторяя процесс, получим вложенные шары $B_k$ исчезающего радиуса.
			\end{itemize}
			Пусть $x_k$ — центры этих шаров. По условию (3) выделим предел $x_k \to x_0$.
			Точка $x_0$ лежит в каком-то открытом $G_{\alpha_0}$. Значит, вокруг $x_0$ есть свободное место (шарик), целиком входящий в $G_{\alpha_0}$.
			
			\textbf{Финал:} С ростом $k$ наши «плохие» шары $B_k$ сжимаются к точке $x_0$ и неизбежно «проваливаются» в это свободное место. Значит, шар $B_k$ накрыт всего **одним** множеством $G_{\alpha_0}$. Противоречие с тем, что он «плохой».
		\end{enumerate}
		
		\item[\textbf{(3) $\Rightarrow$ (4)}] Пусть $M$ бесконечно. Выберем различные $x_n \in M$. По условию выделим сходящуюся подпоследовательность $x_{n_k} \to x_0$. Точка $x_0$ — предельная для $M$.
		
		\item[\textbf{(4) $\Rightarrow$ (3)}] Пусть $\{x_n\}$ — последовательность.
		\begin{itemize}
			\item Если множество значений конечно: выберем постоянную (сходящуюся) подпоследовательность.
			\item Если бесконечно: оно имеет предельную точку $x_0$. Выберем подпоследовательность, сходящуюся к $x_0$.
		\end{itemize}
	\end{itemize}
\end{proof}

\begin{remark}
	Условию компактности пространства $X$ также эквивалентно такое условие: любая непрерывная функция $f : X \to \mathbb{R}$ является ограниченной на $X$, то есть выполнено включение $C(X, \mathbb{R}) \subset B(X, \mathbb{R})$.
\end{remark}

\begin{theorem}[Кантора]\label{thm3.3}
	Пусть $X$ — компактное МП, и функция $f: X \to \mathbb{R}$ непрерывна на $X$. Тогда $f$ равномерно непрерывна на $X$.
\end{theorem}

\begin{definition}[Предкомпакт]
	Пусть $X$ — метрическое пространство. Множество $M \subset X$ называется \textbf{предкомпактом}, если множество $\overline M$ компактно как подпространство в $X$.
\end{definition}
